\documentclass[book.tex]{subfiles}

\begin{document}

\chapter{Electrons in Metals}

\label{chap:electrons_in_metals}
\noindent\rule{11cm}{0.4pt} \\
\textbf{Prerequisites:}  \autoref{chap:ideal_gas}  \\
\textbf{Difficulty Level:} ** \\
\noindent\rule{11cm}{0.4pt}
\vspace{5mm}


Electrons are fundamental elementary particles that are fermions. With this information, together with a correct quantum-mechanical model for the energy levels of electrons, we can derive the thermodynamic properties of electrons present inside metals. These electrons can modeled as an ideal gas of fermions with high densities, assuring the occupation of many single particle energy levels due to the Pauli exclusion principle, and making the interaction between particles negligible.

In this chapter we will develop the thermodynamic description of electrons in metals.


\section{Fermi-Dirac statistics for conducting electrons in a metal}

Electrons in a metal can be modeled as a gas of non-interacting
fermions considering the following assumptions:

\begin{itemize}
\item Electrons in a metal have high densities (many atoms per unit volume each contribute to the conducting electrons).
\item Since no two electrons (fermions) can exist in the same state, a high density system fills many of the single particle energy
levels
\item The lowest unoccupied state will have a kinetic energy much larger than $k_BT$, thus thermal excitations result in energetics
with large kinetic energy and comparatively negligible potential
energy of interaction.
\item The large kinetic energy associated with these electrons results in large conductivity of electrons in metals
\end{itemize}

Using the results for the thermodynamic behavior of non-interacting fermions, we want to find the behavior of an ideal gas of electrons in a metal. Based on the previous assumptions, electrons in a metal act as non-interacting particles with quantized energies given by:
\begin{align}
E_{\vec{n}} = \frac{h^2}{8mL^2}\left(n_x^2+n_y^2+n_z^2\right) \hspace{1.0cm} \text{(translational modes)} \notag
\end{align}

\noindent which are the energy levels for a particle in a box for a particles with mass $m$.%$m=9.10938\times 10^{-31} kg$.

The average number of electrons in the $\vec{n}$ state is:
\begin{align}\label{Eq7}
\langle n_{\vec{n}} \rangle = \frac{1}{e^{\beta (E_{\vec{n}} - \mu)} + 1} = f(E_{\vec{n}})
\end{align}

\noindent where we have defined the Fermi function 
\begin{align}
f(\varepsilon) = \frac{1}{e^{\beta(\varepsilon - \mu)} + 1}
\end{align}.

The total number of electrons is given by:
\begin{align}
\langle N \rangle = 2 \sum_{n_x = 1}^\infty \sum_{n_y = 1}^\infty \sum_{n_z = 1}^\infty \frac{1}{e^{\beta (E_{\vec{n}} - \mu)} + 1} =  2 \sum_{n_x = 1}^\infty \sum_{n_y = 1}^\infty \sum_{n_z = 1}^\infty f(E_{\vec{n}})
\end{align}

\noindent where we include a factor of 2 since the electrons can exist in spin-up and a spin-down states.

For sufficiently large $V$, the spectrum of translational wavemodes is effectively a continuum. Therefore, we can convert this summation to an integral over $\vec{n}$, resulting in:
\begin{align}
\langle N \rangle &= 2 \int_0^\infty dn_x \int_0^\infty dn_y \int_0^\infty dn_z f(E_{\vec{n}})\\
&= 2 \int_0^\infty dk_x \frac{L}{\pi} \int_0^\infty dk_y \frac{L}{\pi} \int_0^\infty dk_z \frac{L}{\pi} f[\varepsilon (\vec{k})] \notag\\
&= 2 \frac{L^3}{\pi^3} \frac{1}{2^3} \int_{-\infty}^\infty dk_x \int_{-\infty}^\infty dk_y \int_{-\infty}^\infty dk_z f[\varepsilon (\vec{k})] \\ 
&= \frac{2V}{(2\pi)^3} \int d\vec{k} F[\varepsilon (\vec{k})] 
\end{align}

\noindent where we have used a coordinate change from $\vec{n}$ to $\vec{k} = \frac{\pi}{L}\vec{n}$, and the energy is now written as \begin{align}
\varepsilon(\vec{k}) = \frac{\hslash^2 k^2}{2m}
\end{align}

Define the chemical potential at $T =0$ to be 
\begin{align} \mu(T=0)=\varepsilon_F\end{align} This is referred to as the Fermi energy. Consider the form of the Fermi function at $T=0$, 
\begin{align}
f(\varepsilon) = \begin{cases}
1 & \textrm{ for }\varepsilon \le \varepsilon_F \\ 
0 & \textrm{ for }\varepsilon > \varepsilon_F
\end{cases}
\end{align}
Define the Fermi momentum $p_F$ according to \begin{align}
\varepsilon_F &= \frac{p^2_F}{2m} = \frac{\hslash^2 k_F^2}{2m}
\end{align}
Thus, at $T = 0$, $\langle N \rangle$ is found to be:
\begin{align} \label{AveN_electrons}
\langle N \rangle = \frac{2V}{(2\pi)^3} \frac{4}{3}\pi k_F^3 = \frac{8}{3}\pi\frac{V(2m)^{3/2}}{h^3}\varepsilon_F^{3/2}
\end{align}
\noindent where the fermi energy $\varepsilon_F$ is given explicitly by:
\begin{align} \label{FermiEnergy}
\varepsilon_F = \left(\frac{3\rho}{8\pi}\right)^{2/3} \frac{h^2}{2m}
\end{align}

A typical metal (\textit{Cu}) has a mass density of $9g/cm^3$. Assuming each atom donates a single electron to the conducting electron gas, this density has a Fermi temperature:
\begin{align}
\theta_F = \frac{\varepsilon_F}{k_B} \approx 80,000K
\end{align}

\noindent this verifies that the Fermi energy $\varepsilon_F$ is sufficiently large to make the ideal gas approximation valid at room temperature. 
At room temperature, only states with energy very near $\varepsilon \approx \varepsilon_F$ will be affected by thermal energy $k_BT$. The spread in the distribution is approximately $2k_BT$ (Fig.~\ref{fig1}).

\begin{marginfigure}
\includegraphics[width=\linewidth]{figures/Physics/statistical_mechanics/electrons_in_metals/Fig1.png}
\caption{The Fermi function $F(\varepsilon) = \left[e^{\beta(\varepsilon - \mu)} +1\right]^{-1}$}
\label{fig1}
\end{marginfigure}

The pressure is found using the relationship for Fermi particles
\begin{equation}\label{Eq9}
pV = 2k_BT\sum_\alpha \log \left(1 + e^{-\beta \varepsilon_\alpha + \beta \mu} \right)
\end{equation}

Following a similar derivation as before, we write
\begin{align}\label{Eq10}
pV &= \frac{2V}{(2\pi)^3} \int d\vec{k} \log \{ 1 + e^{-\beta [ \varepsilon (\vec{k}) - \mu]} \} \\
&= \frac{4\pi V (2m)^{2/3}}{h^3} \int_0^\infty d\varepsilon \varepsilon^{1/2} \log \left[ 1 + e^{-\beta ( \varepsilon - \mu)} \right]
\end{align}

\noindent where we have used $\varepsilon = \frac{\hslash^2 k^2}{2m}$.

In the limit $T \rightarrow 0$ or $\beta \rightarrow \infty$, we have 
\begin{align} 
1 + e^{-\beta ( \varepsilon - \mu)} \rightarrow  e^{\beta ( \varepsilon_F - \varepsilon)}
\end{align}
for $\varepsilon < \varepsilon_F$ and 
\begin{align} 1 + e^{-\beta ( \varepsilon - \mu)} \rightarrow 1
\end{align}
for $\varepsilon > \varepsilon_F$. Therefore, the pressure is written as:

\begin{align}
pV &= \frac{4\pi V (2m)^{3/2}}{h^3} \int_0^{\varepsilon_F} d\varepsilon \varepsilon^{1/2} (\varepsilon_F - \varepsilon)\\
&= \frac{16\pi V (2m)^{2/3}}{15h^3} \varepsilon_F^{5/2} 
\end{align}

This pressure at $T = 0$ is approximately $10^6 \textrm{ atm}$. This large pressure plays an important role in halting the collapse of a star
(white dwarf) because this enormous pressure offsets the gravitational forces that otherwise drive the collapse.

The average energy $\langle E \rangle$ is found using:
\begin{align}\label{Eq12}
\langle E \rangle = 2\sum_\alpha \varepsilon_\alpha \frac{e^{-\beta \varepsilon_\alpha + \beta \mu}}{1 + e^{-\beta \varepsilon_\alpha + \beta \mu}}
\end{align}

Following a similar derivation as before, we write
\begin{align}
\langle E \rangle &= \frac{2V}{(2\pi)^3} \int d\vec{k} \varepsilon(\vec{k}) \frac{1}{e^{\beta [ \varepsilon(\vec{k}) - \mu]}+1} \notag\\
&= \frac{4\pi V (2m)^{2/3}}{h^3} \int_0^\infty d\varepsilon \varepsilon^{3/2} \frac{1}{e^{\beta [ \varepsilon - \mu]}+1} \notag\\
&= \frac{4\pi V (2m)^{2/3}}{h^3} \int_0^\infty d\varepsilon \varepsilon^{3/2} F(\varepsilon) \notag
\end{align}

\noindent where we have used $\varepsilon = \frac{\hslash^2 k^2}{2m}$.

\begin{marginfigure}
\includegraphics[width=\linewidth]{figures/Physics/statistical_mechanics/electrons_in_metals/Fig2.png}
\caption{The derivative of the Fermi function: \\ \\
$\frac{df(\varepsilon)}{d\varepsilon} = \beta e^{\beta (\varepsilon - \mu)}{[e^{\beta (\varepsilon - \mu)} + 1]}^{-2}$
}
\label{fig2}
\end{marginfigure}

We apply integration by parts to this equation and use (\ref{AveN_electrons}) to get:
\begin{align}
\langle E \rangle = - \frac{3\langle N \rangle}{5\varepsilon_F^{3/2}} \int_0^{\infty} d\varepsilon \varepsilon^{5/2} \frac{df}{d\varepsilon}
\end{align}

\noindent where 
\begin{align}
\frac{df}{d\varepsilon} &= -\frac{\beta e^{\beta (\varepsilon - \mu)}}{[e^{\beta (\varepsilon - \mu)} + 1]^2}
\end{align}

In the limit $T\rightarrow0$, the function $-\frac{df}{d\varepsilon}$ becomes peaked near $\varepsilon \approx \varepsilon_F$  (Fig.~\ref{fig2}), and we can effectively expand the integrand near $\varepsilon = \varepsilon_F$  to get:
\begin{align}
\varepsilon^{5/2} \approx \varepsilon_F^{5/2} + \frac{5}{2}\varepsilon_F^{3/2}(\varepsilon - \varepsilon_F) + \frac{15}{8}\varepsilon_F^{1/2}(\varepsilon - \varepsilon_F)^2 + ...
\end{align}
 
%\begin{marginfigure}
%\includegraphics[width=\linewidth]{figures/Physics/statistical_mechanics/electrons_in_metals/Fig3.png}
%\caption{The heat capacity $C_V$ of $Ti_3SiC_2$ exhibits the temperature dependence $C_V=\gamma_1T+\gamma_3T^3+...$ in the low-temperature limit \citep{Ho1999}}
%\label{fig3}
%\end{marginfigure}

Since $\frac{df(\varepsilon)}{d\varepsilon}$ is even about $\varepsilon_F$ in this limit, the odd-order terms will integrate to zero, leaving only the even terms. Thus, the final form of the average energy is going to be:
\begin{align}
\langle E \rangle = \langle N \rangle \varepsilon_F \left[A +B\left(\frac{T}{\theta_F}\right)^2 +...\right]
\end{align}

A precise calculation of this low-temperature expansion gives:
\begin{align}\
\langle E \rangle = \frac{3}{5}\langle N \rangle \varepsilon_F \left[1 +\frac{5\pi^2}{12}\left(\frac{T}{\theta_F}\right)^2 +...\right]
\end{align}

Therefore, the heat capacity for a metal in the limit of small temperature is given by

\begin{align}
C_V = \frac{3}{5}\langle N \rangle \varepsilon_F \frac{5\pi^2}{12}2\frac{T}{\theta_F^2} = \frac{\pi^2}{2}\langle N \rangle k_B\frac{T}{\theta_F}
\end{align}

\noindent giving a linear temperature dependence in the small-$T$ limit.

The limiting behavior in the small-$T$ limit suggests that a plot of $C_V/T$ approaches a constant as $T\rightarrow 0$. This proves to be the behavior of the heat capacity for metals in the limit of small $T$. As temperature increases, the fluctuations in the metal nuclei also contribute to the heat capacity (Fig.~\ref{fig3}).

\end{document}